% Generated by export_prompts.py -- do not edit by hand.
\begin{enumerate}[label=\textbf{P\arabic*.},leftmargin=2.6em,itemsep=4pt]
\item {\footnotesize\color{black!60}2026-11-02}
\begin{lstlisting}[style=prompt]
Find out how you can be a TA for me that guides me to the correct answer but don't tell me exactly what the answer is based on @hw4-sp26-cs633.pdf @hw4-data/ @hw4_ske.py
\end{lstlisting}
\item {\footnotesize\color{black!60}2026-11-02}
\begin{lstlisting}[style=prompt]
prepare the virtual environment using uv
\end{lstlisting}
\item {\footnotesize\color{black!60}2026-11-03}
\begin{lstlisting}[style=prompt]
what methods are available to me by inheriting Dataset?
\end{lstlisting}
\item {\footnotesize\color{black!60}2026-11-03}
\begin{lstlisting}[style=prompt]
What do you mean by map-style? Also it doesn't seem right to load all images into ram. What I was thinking is maybe using a map that uses ints as a way to index into the map, and then return a tuple of (path, label)
\end{lstlisting}
\item {\footnotesize\color{black!60}2026-11-03}
\begin{lstlisting}[style=prompt]
how do I use with open(path, r) as f:?
\end{lstlisting}
\item {\footnotesize\color{black!60}2026-11-03}
\begin{lstlisting}[style=prompt]
theres a file called "ClassName.txt" that contains the names of the class. I'm using that as our index for integer labels
\end{lstlisting}
\item {\footnotesize\color{black!60}2026-11-03}
\begin{lstlisting}[style=prompt]
i want to do it from the folder structure directly, but I don't remember the exact syntax. I know it is in the os library
\end{lstlisting}
\item {\footnotesize\color{black!60}2026-11-03}
\begin{lstlisting}[style=prompt]
~~~~~~~~~~~~~~~~~~~~~^^
  File "/Users/arkan/Documents/School/Machine-Learning/HW/4/hw4_ske.py", line 29, in populate_dataset
    sorted_dirs = sorted(os.listdir(dir_path))
                         ~~~~~~~~~~^^^^^^^^^^
NotADirectoryError: [Errno 20] Not a directory: 'hw4-data/.DS_Store'
\end{lstlisting}
\item {\footnotesize\color{black!60}2026-11-03}
\begin{lstlisting}[style=prompt]
~~~~~~~~~~~~~~~~~~~~~^^
[... 16 lines omitted ...]
 Why is it not working?
\end{lstlisting}
\item {\footnotesize\color{black!60}2026-11-03}
\begin{lstlisting}[style=prompt]
E 46  ?   ?   ?   ?   self.data.append((os.path.abspath(os.path.join(dir_path, f)), label))
[... 9 lines omitted ...]
?        "__fspath__" is not present [reportArgumentType]                                                                ?
\end{lstlisting}
\item {\footnotesize\color{black!60}2026-11-03}
\begin{lstlisting}[style=prompt]
for f in filenames:
                label = np.zeros(n)
                label[i] = 1.0
                self.data.append((os.path.abspath(os.path.join(dir_path, f)), label))
\end{lstlisting}
\item {\footnotesize\color{black!60}2026-11-03}
\begin{lstlisting}[style=prompt]
filenames comes from os.walk, also it says label tensor, so that implies multiple values
\end{lstlisting}
\item {\footnotesize\color{black!60}2026-11-03}
\begin{lstlisting}[style=prompt]
n = len(labels)
[... 8 lines omitted ...]
                self.data.append((os.path.abspath(os.path.join(dir_path, f)), label))
\end{lstlisting}
\item {\footnotesize\color{black!60}2026-11-04}
\begin{lstlisting}[style=prompt]
how do I use PIL.Image and convert it into a numpy array?
\end{lstlisting}
\item {\footnotesize\color{black!60}2026-11-04}
\begin{lstlisting}[style=prompt]
back to the crossentropyloss discussion. my understanding is that the classes should either be yes or no. in a multiclass classfier, we should use a one-hot encoding
\end{lstlisting}
\item {\footnotesize\color{black!60}2026-11-04}
\begin{lstlisting}[style=prompt]
populate .gitignore for me
\end{lstlisting}
\item {\footnotesize\color{black!60}2026-11-04}
\begin{lstlisting}[style=prompt]
when i tested the dataset, it says that Test Failed: 314 != 3 : Images should have 3 channels (RGB). why is that?
\end{lstlisting}
\item {\footnotesize\color{black!60}2026-11-04}
\begin{lstlisting}[style=prompt]
i'm confused. I don't get why when I'm trying to print the kwargs, I get
[... 8 lines omitted ...]
 What is really the input of the calculate_mean_std function?
\end{lstlisting}
\item {\footnotesize\color{black!60}2026-11-04}
\begin{lstlisting}[style=prompt]
calculating the mean of the channels seems pretty straight forward, but finding the std does not. for mean, I can imagine summing up all the values and then dividing it by the number of pixels counted for all images, but i don't know if that approach works for standard deviation. guide me to the right answer
\end{lstlisting}
\item {\footnotesize\color{black!60}2026-11-04}
\begin{lstlisting}[style=prompt]
1. they are pretty much the same as in the same number of computation happens. 2. i believe if i convert all the numbers from [0-255] to [0-1] i won't get that issue 3. I want to use the kwargs functionality for calculate. if kwargs is empty, return the hard coded values. if not, find out what the mean_std is based on the given data. that way there wouldn't be any data leak. i would just need to figure out how to do train/val split in dataloader
\end{lstlisting}
\item {\footnotesize\color{black!60}2026-11-04}
\begin{lstlisting}[style=prompt]
i'll create a new function that takes in dataloader as a parameter. still though, that doesn't answer how train/val split works. does doing so spawn a new dataloader? what's the mechanism behind that
\end{lstlisting}
\item {\footnotesize\color{black!60}2026-11-04}
\begin{lstlisting}[style=prompt]
this means that full_dataset inherits DataSet right?
\end{lstlisting}
\item {\footnotesize\color{black!60}2026-11-04}
\begin{lstlisting}[style=prompt]
how do I use dataloader? is it the same as getitem? how do I iterate over it?
\end{lstlisting}
\item {\footnotesize\color{black!60}2026-11-04}
\begin{lstlisting}[style=prompt]
can I do vectorized numpy operations on images? or do I have to convert it. also, if i'm given batches, do I sum up a channel of multiple images?
\end{lstlisting}
\item {\footnotesize\color{black!60}2026-11-04}
\begin{lstlisting}[style=prompt]
can I assert that batch_size == 1?
\end{lstlisting}
\item {\footnotesize\color{black!60}2026-11-04}
\begin{lstlisting}[style=prompt]
i want to use dataloader so that I have the potential of avoiding data leakage, I want to use offline_mean_std for full dataset, train, and val
\end{lstlisting}
\item {\footnotesize\color{black!60}2026-11-04}
\begin{lstlisting}[style=prompt]
offline_mean_std(dataloader), also the full dataset should be the one that get's hardcoded, that means that inevitably i'll get data leakage
\end{lstlisting}
\item {\footnotesize\color{black!60}2026-11-04}
\begin{lstlisting}[style=prompt]
calculate_mean_std is going to be the function i turn in for the assignment. I won't be using it otherwise
\end{lstlisting}
\item {\footnotesize\color{black!60}2026-11-04}
\begin{lstlisting}[style=prompt]
computes from offline_mean_std
\end{lstlisting}
\item {\footnotesize\color{black!60}2026-11-04}
\begin{lstlisting}[style=prompt]
as I work through this, document it in a file in this directory. and make sure it's easily imported for latex documents
\end{lstlisting}
\item {\footnotesize\color{black!60}2026-11-04}
\begin{lstlisting}[style=prompt]
can I assert a paramter of dataloader in offline_mean_std? how would I do that?
\end{lstlisting}
\item {\footnotesize\color{black!60}2026-11-04}
\begin{lstlisting}[style=prompt]
also, do it proactively in natural checkpoints
\end{lstlisting}
\item {\footnotesize\color{black!60}2026-11-04}
\begin{lstlisting}[style=prompt]
now that batch size is one, I should only get a 3-d array instead of 4. how do I sum now?
\end{lstlisting}
\item {\footnotesize\color{black!60}2026-11-04}
\begin{lstlisting}[style=prompt]
okay, and how do I square the values of certain axises?
\end{lstlisting}
\item {\footnotesize\color{black!60}2026-11-04}
\begin{lstlisting}[style=prompt]
wait, does sum(axis=(0,2,3)) return 3 numbers?
\end{lstlisting}
\item {\footnotesize\color{black!60}2026-11-04}
\begin{lstlisting}[style=prompt]
how do I elegantly get the values of H and W, multiply them, and keep a running total of num pixels?
\end{lstlisting}
\item {\footnotesize\color{black!60}2026-11-04}
\begin{lstlisting}[style=prompt]
can I unpack a np 1d array of size 3 into 3 different variables using a, b, c = arr?
\end{lstlisting}
\item {\footnotesize\color{black!60}2026-11-04}
\begin{lstlisting}[style=prompt]
review my offline_mean_std, does it look good?
\end{lstlisting}
\item {\footnotesize\color{black!60}2026-11-04}
\begin{lstlisting}[style=prompt]
def offline_mean_std(dataloader: DataLoader):
[... 22 lines omitted ...]
    return [mean_r, mean_g, mean_b], [std_r, std_g, std_b]
\end{lstlisting}
\item {\footnotesize\color{black!60}2026-11-04}
\begin{lstlisting}[style=prompt]
does torch tensor have the same capabilities as numpy?
\end{lstlisting}
\item {\footnotesize\color{black!60}2026-11-04}
\begin{lstlisting}[style=prompt]
does tracking gradients mess up with training at all here if I compute in torch instad of nunmpy?
\end{lstlisting}
\item {\footnotesize\color{black!60}2026-11-04}
\begin{lstlisting}[style=prompt]
how do I make torch.zeros(3) into a float64 value
\end{lstlisting}
\item {\footnotesize\color{black!60}2026-11-04}
\begin{lstlisting}[style=prompt]
def offline_mean_std(dataloader: DataLoader):
[... 22 lines omitted ...]
    return [mean_r, mean_g, mean_b], [std_r, std_g, std_b]
\end{lstlisting}
\item {\footnotesize\color{black!60}2026-11-04}
\begin{lstlisting}[style=prompt]
could I do float64 in the dataset? in __getitem__?
\end{lstlisting}
\item {\footnotesize\color{black!60}2026-11-04}
\begin{lstlisting}[style=prompt]
so how would I do that in offline_mean_std before it gets squared?
\end{lstlisting}
\item {\footnotesize\color{black!60}2026-11-04}
\begin{lstlisting}[style=prompt]
let's start working on the cnn. I could probably just define all the layers in the forward pass right?
\end{lstlisting}
\item {\footnotesize\color{black!60}2026-11-04}
\begin{lstlisting}[style=prompt]
I want to use Conv3d since we are dealing with rgb images. I also don't know the parameters for Conv3d
\end{lstlisting}
\item {\footnotesize\color{black!60}2026-11-04}
\begin{lstlisting}[style=prompt]
out_channels return the number of channels similar to in_channels. kernel size refers to how big the window is, if its 3x3 or something else
\end{lstlisting}
\item {\footnotesize\color{black!60}2026-11-04}
\begin{lstlisting}[style=prompt]
i'm not sure if how out_channels will help me, and if I should make it 16. the second layer in_channel should always match the first layers out_channel. stride does change the output spatial size. i don't think i should add any padding aside from 0 padding when the kernel is out of the frame
\end{lstlisting}
\item {\footnotesize\color{black!60}2026-11-04}
\begin{lstlisting}[style=prompt]
i want to start off with 2 conv layers, a max pooling layer, one fc, and then one softmax. which classes should i use?
\end{lstlisting}
\item {\footnotesize\color{black!60}2026-11-04}
\begin{lstlisting}[style=prompt]
can I reuse the same nn.ReLU object multiple times?
\end{lstlisting}
\item {\footnotesize\color{black!60}2026-11-04}
\begin{lstlisting}[style=prompt]
what should the kernel_size be for 224x224 images
\end{lstlisting}
\item {\footnotesize\color{black!60}2026-11-04}
\begin{lstlisting}[style=prompt]
also should the number of channels reduce per layer?
\end{lstlisting}
\item {\footnotesize\color{black!60}2026-11-04}
\begin{lstlisting}[style=prompt]
what does max pool per layer do?
\end{lstlisting}
\item {\footnotesize\color{black!60}2026-11-04}
\begin{lstlisting}[style=prompt]
I'm pretty sure i do not need to calculate all that. there must be a simpler way to just get the number of inputs for linear without all of that
\end{lstlisting}
\item {\footnotesize\color{black!60}2026-11-04}
\begin{lstlisting}[style=prompt]
is out_features for the fc layer num_classes?
\end{lstlisting}
\item {\footnotesize\color{black!60}2026-11-04}
\begin{lstlisting}[style=prompt]
def __init__(self, num_classes=10):
[... 48 lines omitted ...]
        return self.fc(x)
\end{lstlisting}
\item {\footnotesize\color{black!60}2026-11-04}
\begin{lstlisting}[style=prompt]
what would be the best padding per layer? also how do I use flatten?
\end{lstlisting}
\item {\footnotesize\color{black!60}2026-11-04}
\begin{lstlisting}[style=prompt]
lets use maxpool2d to be the thing that controls downsampling, so that means it'll be after every conv1 + relu layer
\end{lstlisting}
\item {\footnotesize\color{black!60}2026-11-04}
\begin{lstlisting}[style=prompt]
class CNN(nn.Module):
[... 54 lines omitted ...]
 i want to do this
\end{lstlisting}
\item {\footnotesize\color{black!60}2026-11-04}
\begin{lstlisting}[style=prompt]
def forward(self, x):
[... 23 lines omitted ...]
        return self.fc(x)
\end{lstlisting}
\item {\footnotesize\color{black!60}2026-11-04}
\begin{lstlisting}[style=prompt]
def calculate_fc_input():
            dummy = torch.zeros(1, 3, 224, 224)
            out = self.pool(self.conv3(self.pool(self.conv2(self.pool(self.conv1(dummy))))))
            return out.shape
\end{lstlisting}
\item {\footnotesize\color{black!60}2026-11-04}
\begin{lstlisting}[style=prompt]
class CNN(nn.Module):
[... 53 lines omitted ...]
        return self.fc(x)
\end{lstlisting}
\item {\footnotesize\color{black!60}2026-11-04}
\begin{lstlisting}[style=prompt]
let's work on normalization
\end{lstlisting}
\item {\footnotesize\color{black!60}2026-11-04}
\begin{lstlisting}[style=prompt]
I'm going to use this
[... 9 lines omitted ...]
 I also will create a normalize global function
\end{lstlisting}
\item {\footnotesize\color{black!60}2026-11-04}
\begin{lstlisting}[style=prompt]
I will be using only the train mean std. I will hardcode numbers separately. I didn't even think of using combiging ToTensor and normalize in on transform function. Does Normalize(mean, std) take into list/numpy arrays?
\end{lstlisting}
\item {\footnotesize\color{black!60}2026-11-04}
\begin{lstlisting}[style=prompt]
calculate_mean_std contains the mean std for the full dataset, i don't want data leakage
\end{lstlisting}
\item {\footnotesize\color{black!60}2026-11-04}
\begin{lstlisting}[style=prompt]
What I'll do most likely is use the train mean std for local training. How do I build the transform compose function thing
\end{lstlisting}
\item {\footnotesize\color{black!60}2026-11-04}
\begin{lstlisting}[style=prompt]
how do I use a transform?
\end{lstlisting}
\item {\footnotesize\color{black!60}2026-11-04}
\begin{lstlisting}[style=prompt]
it seems like the mean std stats for all 3 aren't that different after all. i'll just use calculate_mean_std
\end{lstlisting}
\item {\footnotesize\color{black!60}2026-11-04}
\begin{lstlisting}[style=prompt]
is there another transform that clips images to 224x224?
\end{lstlisting}
\item {\footnotesize\color{black!60}2026-11-04}
\begin{lstlisting}[style=prompt]
Right after normalize
\end{lstlisting}
\item {\footnotesize\color{black!60}2026-11-04}
\begin{lstlisting}[style=prompt]
can you do resize on a pil even though its from pytorch?
\end{lstlisting}
\item {\footnotesize\color{black!60}2026-11-04}
\begin{lstlisting}[style=prompt]
for training, which kwargs should I be aware of? I know of different optimizers, epochs, etc. I'm going to employ a custom grid search mechanism to find hyperparameters
\end{lstlisting}
\item {\footnotesize\color{black!60}2026-11-04}
\begin{lstlisting}[style=prompt]
i'm using a mac m5. what hsould be the device?
\end{lstlisting}
\item {\footnotesize\color{black!60}2026-11-04}
\begin{lstlisting}[style=prompt]
does using certain devices change the structure of models.pt?
\end{lstlisting}
\item {\footnotesize\color{black!60}2026-11-04}
\begin{lstlisting}[style=prompt]
what is the pseudocode for train and test?
\end{lstlisting}
\item {\footnotesize\color{black!60}2026-11-04}
\begin{lstlisting}[style=prompt]
for the hyperparameters, give me a clean list of which ones I should tune and the range/options of each one
\end{lstlisting}
\item {\footnotesize\color{black!60}2026-11-04}
\begin{lstlisting}[style=prompt]
are the name of the hyperparameters the exact keyword that we pass in? how do I createoptimizer and loss function and scheduler?
\end{lstlisting}
\item {\footnotesize\color{black!60}2026-11-04}
\begin{lstlisting}[style=prompt]
how do I use test function as the val accuracy?
\end{lstlisting}
\item {\footnotesize\color{black!60}2026-11-04}
\begin{lstlisting}[style=prompt]
is model.to(device=device) idempotent?
\end{lstlisting}
\item {\footnotesize\color{black!60}2026-11-04}
\begin{lstlisting}[style=prompt]
how to disable grad tracking, and reneable it
\end{lstlisting}
\item {\footnotesize\color{black!60}2026-11-04}
\begin{lstlisting}[style=prompt]
how do I increment total and correct?
\end{lstlisting}
\item {\footnotesize\color{black!60}2026-11-04}
\begin{lstlisting}[style=prompt]
how should I initialize optim.Adam
\end{lstlisting}
\item {\footnotesize\color{black!60}2026-11-04}
\begin{lstlisting}[style=prompt]
how to zero out optimzer acculated gradients
\end{lstlisting}
\item {\footnotesize\color{black!60}2026-11-04}
\begin{lstlisting}[style=prompt]
based on my grid_search, how do I create the scheduler?
\end{lstlisting}
\item {\footnotesize\color{black!60}2026-11-04}
\begin{lstlisting}[style=prompt]
how do I update running_loss in train?
\end{lstlisting}
\item {\footnotesize\color{black!60}2026-11-04}
\begin{lstlisting}[style=prompt]
this means I should keep a running total of images seen as well so that I can get a good mean
\end{lstlisting}
\item {\footnotesize\color{black!60}2026-11-04}
\begin{lstlisting}[style=prompt]
how do I use rich.progress to log the loops and everything
\end{lstlisting}
\item {\footnotesize\color{black!60}2026-11-04}
\begin{lstlisting}[style=prompt]
who it be batch_task advance=1 or advance=batch_size?
\end{lstlisting}
\item {\footnotesize\color{black!60}2026-11-04}
\begin{lstlisting}[style=prompt]
what's the default for DataLoader classes?
\end{lstlisting}
\item {\footnotesize\color{black!60}2026-11-04}
\begin{lstlisting}[style=prompt]
how do I report running_loss for epoch, curr_val_accuracy, using progress
\end{lstlisting}
\item {\footnotesize\color{black!60}2026-11-04}
\begin{lstlisting}[style=prompt]
how do I log the fact that val_accuracy for a certain epoch is the best and that I want to inform the user that I saved it
\end{lstlisting}
\item {\footnotesize\color{black!60}2026-11-04}
\begin{lstlisting}[style=prompt]
I don't want to save model.pt into a new file. I want to keep it as a variable for now. is that possible?
\end{lstlisting}
\item {\footnotesize\color{black!60}2026-11-04}
\begin{lstlisting}[style=prompt]
is copy part of standard python library?
\end{lstlisting}
\item {\footnotesize\color{black!60}2026-11-04}
\begin{lstlisting}[style=prompt]
def train(model: CNN, train_loader: DataLoader, val_loader: DataLoader, **kwargs):
[... 44 lines omitted ...]
        return best_model_weights, best_val_accuracy
\end{lstlisting}
\item {\footnotesize\color{black!60}2026-11-04}
\begin{lstlisting}[style=prompt]
I want to add a task in progress that says what hyperparamters I'm using, how do I do that?
\end{lstlisting}
\item {\footnotesize\color{black!60}2026-11-04}
\begin{lstlisting}[style=prompt]
read through @hw4_ske.py and review for any errors
\end{lstlisting}
\item {\footnotesize\color{black!60}2026-11-04}
\begin{lstlisting}[style=prompt]
running_loss = 0
            total_images = 0
            for images, labels in train_loader:
                images = images.to(device)
                labels = labels.to(device)
\end{lstlisting}
\item {\footnotesize\color{black!60}2026-11-05}
\begin{lstlisting}[style=prompt]
review @hw4_ske.py
\end{lstlisting}
\item {\footnotesize\color{black!60}2026-11-05}
\begin{lstlisting}[style=prompt]
progress.update(
[... 5 lines omitted ...]
 this line does not change description at all
\end{lstlisting}
\item {\footnotesize\color{black!60}2026-11-05}
\begin{lstlisting}[style=prompt]
with Progress(
[... 43 lines omitted ...]
    return best_model, best_val_accuracy
\end{lstlisting}
\item {\footnotesize\color{black!60}2026-11-05}
\begin{lstlisting}[style=prompt]
for some reason, loss is stuck at 1.3322 and the val_acc is not going up or down
\end{lstlisting}
\item {\footnotesize\color{black!60}2026-11-05}
\begin{lstlisting}[style=prompt]
why don't you check all that for me and tell me what is going on?
\end{lstlisting}
\item {\footnotesize\color{black!60}2026-11-05}
\begin{lstlisting}[style=prompt]
find out what is structurally wrong with my code. My loss is not going down
\end{lstlisting}
\item {\footnotesize\color{black!60}2026-11-05}
\begin{lstlisting}[style=prompt]
could I resue the same batchnorm layer?
\end{lstlisting}
\item {\footnotesize\color{black!60}2026-11-05}
\begin{lstlisting}[style=prompt]
yes
\end{lstlisting}
\item {\footnotesize\color{black!60}2026-11-05}
\begin{lstlisting}[style=prompt]
I don't see those changes
\end{lstlisting}
\item {\footnotesize\color{black!60}2026-11-05}
\begin{lstlisting}[style=prompt]
my val accuracy never went higher than 51.9%. what else can i fix?
\end{lstlisting}
\item {\footnotesize\color{black!60}2026-11-05}
\begin{lstlisting}[style=prompt]
update the @hw4_ske.py with the new changes
\end{lstlisting}
\item {\footnotesize\color{black!60}2026-11-05}
\begin{lstlisting}[style=prompt]
is the code running in mps or cpu?
\end{lstlisting}
\item {\footnotesize\color{black!60}2026-11-02}
\begin{lstlisting}[style=prompt]
Find out how you can be a TA for me that guides me to the correct answer but don't tell me exactly what the answer is based on @hw4-sp26-cs633.pdf @hw4-data/ @hw4_ske.py
\end{lstlisting}
\item {\footnotesize\color{black!60}2026-11-02}
\begin{lstlisting}[style=prompt]
prepare the virtual environment using uv
\end{lstlisting}
\item {\footnotesize\color{black!60}2026-11-03}
\begin{lstlisting}[style=prompt]
what methods are available to me by inheriting Dataset?
\end{lstlisting}
\item {\footnotesize\color{black!60}2026-11-03}
\begin{lstlisting}[style=prompt]
What do you mean by map-style? Also it doesn't seem right to load all images into ram. What I was thinking is maybe using a map that uses ints as a way to index into the map, and then return a tuple of (path, label)
\end{lstlisting}
\item {\footnotesize\color{black!60}2026-11-03}
\begin{lstlisting}[style=prompt]
how do I use with open(path, r) as f:?
\end{lstlisting}
\item {\footnotesize\color{black!60}2026-11-03}
\begin{lstlisting}[style=prompt]
theres a file called "ClassName.txt" that contains the names of the class. I'm using that as our index for integer labels
\end{lstlisting}
\item {\footnotesize\color{black!60}2026-11-03}
\begin{lstlisting}[style=prompt]
i want to do it from the folder structure directly, but I don't remember the exact syntax. I know it is in the os library
\end{lstlisting}
\item {\footnotesize\color{black!60}2026-11-03}
\begin{lstlisting}[style=prompt]
~~~~~~~~~~~~~~~~~~~~~^^
  File "/Users/arkan/Documents/School/Machine-Learning/HW/4/hw4_ske.py", line 29, in populate_dataset
    sorted_dirs = sorted(os.listdir(dir_path))
                         ~~~~~~~~~~^^^^^^^^^^
NotADirectoryError: [Errno 20] Not a directory: 'hw4-data/.DS_Store'
\end{lstlisting}
\item {\footnotesize\color{black!60}2026-11-03}
\begin{lstlisting}[style=prompt]
~~~~~~~~~~~~~~~~~~~~~^^
[... 16 lines omitted ...]
 Why is it not working?
\end{lstlisting}
\item {\footnotesize\color{black!60}2026-11-03}
\begin{lstlisting}[style=prompt]
E 46  ?   ?   ?   ?   self.data.append((os.path.abspath(os.path.join(dir_path, f)), label))
[... 9 lines omitted ...]
?        "__fspath__" is not present [reportArgumentType]                                                                ?
\end{lstlisting}
\item {\footnotesize\color{black!60}2026-11-03}
\begin{lstlisting}[style=prompt]
for f in filenames:
                label = np.zeros(n)
                label[i] = 1.0
                self.data.append((os.path.abspath(os.path.join(dir_path, f)), label))
\end{lstlisting}
\item {\footnotesize\color{black!60}2026-11-03}
\begin{lstlisting}[style=prompt]
filenames comes from os.walk, also it says label tensor, so that implies multiple values
\end{lstlisting}
\item {\footnotesize\color{black!60}2026-11-03}
\begin{lstlisting}[style=prompt]
n = len(labels)
[... 8 lines omitted ...]
                self.data.append((os.path.abspath(os.path.join(dir_path, f)), label))
\end{lstlisting}
\item {\footnotesize\color{black!60}2026-11-04}
\begin{lstlisting}[style=prompt]
how do I use PIL.Image and convert it into a numpy array?
\end{lstlisting}
\item {\footnotesize\color{black!60}2026-11-04}
\begin{lstlisting}[style=prompt]
back to the crossentropyloss discussion. my understanding is that the classes should either be yes or no. in a multiclass classfier, we should use a one-hot encoding
\end{lstlisting}
\item {\footnotesize\color{black!60}2026-11-04}
\begin{lstlisting}[style=prompt]
populate .gitignore for me
\end{lstlisting}
\item {\footnotesize\color{black!60}2026-11-04}
\begin{lstlisting}[style=prompt]
when i tested the dataset, it says that Test Failed: 314 != 3 : Images should have 3 channels (RGB). why is that?
\end{lstlisting}
\item {\footnotesize\color{black!60}2026-11-04}
\begin{lstlisting}[style=prompt]
i'm confused. I don't get why when I'm trying to print the kwargs, I get
[... 8 lines omitted ...]
 What is really the input of the calculate_mean_std function?
\end{lstlisting}
\item {\footnotesize\color{black!60}2026-11-04}
\begin{lstlisting}[style=prompt]
calculating the mean of the channels seems pretty straight forward, but finding the std does not. for mean, I can imagine summing up all the values and then dividing it by the number of pixels counted for all images, but i don't know if that approach works for standard deviation. guide me to the right answer
\end{lstlisting}
\item {\footnotesize\color{black!60}2026-11-04}
\begin{lstlisting}[style=prompt]
1. they are pretty much the same as in the same number of computation happens. 2. i believe if i convert all the numbers from [0-255] to [0-1] i won't get that issue 3. I want to use the kwargs functionality for calculate. if kwargs is empty, return the hard coded values. if not, find out what the mean_std is based on the given data. that way there wouldn't be any data leak. i would just need to figure out how to do train/val split in dataloader
\end{lstlisting}
\item {\footnotesize\color{black!60}2026-11-04}
\begin{lstlisting}[style=prompt]
i'll create a new function that takes in dataloader as a parameter. still though, that doesn't answer how train/val split works. does doing so spawn a new dataloader? what's the mechanism behind that
\end{lstlisting}
\item {\footnotesize\color{black!60}2026-11-04}
\begin{lstlisting}[style=prompt]
this means that full_dataset inherits DataSet right?
\end{lstlisting}
\item {\footnotesize\color{black!60}2026-11-04}
\begin{lstlisting}[style=prompt]
how do I use dataloader? is it the same as getitem? how do I iterate over it?
\end{lstlisting}
\item {\footnotesize\color{black!60}2026-11-04}
\begin{lstlisting}[style=prompt]
can I do vectorized numpy operations on images? or do I have to convert it. also, if i'm given batches, do I sum up a channel of multiple images?
\end{lstlisting}
\item {\footnotesize\color{black!60}2026-11-04}
\begin{lstlisting}[style=prompt]
can I assert that batch_size == 1?
\end{lstlisting}
\item {\footnotesize\color{black!60}2026-11-04}
\begin{lstlisting}[style=prompt]
i want to use dataloader so that I have the potential of avoiding data leakage, I want to use offline_mean_std for full dataset, train, and val
\end{lstlisting}
\item {\footnotesize\color{black!60}2026-11-04}
\begin{lstlisting}[style=prompt]
offline_mean_std(dataloader), also the full dataset should be the one that get's hardcoded, that means that inevitably i'll get data leakage
\end{lstlisting}
\item {\footnotesize\color{black!60}2026-11-04}
\begin{lstlisting}[style=prompt]
calculate_mean_std is going to be the function i turn in for the assignment. I won't be using it otherwise
\end{lstlisting}
\item {\footnotesize\color{black!60}2026-11-04}
\begin{lstlisting}[style=prompt]
computes from offline_mean_std
\end{lstlisting}
\item {\footnotesize\color{black!60}2026-11-04}
\begin{lstlisting}[style=prompt]
as I work through this, document it in a file in this directory. and make sure it's easily imported for latex documents
\end{lstlisting}
\item {\footnotesize\color{black!60}2026-11-04}
\begin{lstlisting}[style=prompt]
can I assert a paramter of dataloader in offline_mean_std? how would I do that?
\end{lstlisting}
\item {\footnotesize\color{black!60}2026-11-04}
\begin{lstlisting}[style=prompt]
also, do it proactively in natural checkpoints
\end{lstlisting}
\item {\footnotesize\color{black!60}2026-11-04}
\begin{lstlisting}[style=prompt]
now that batch size is one, I should only get a 3-d array instead of 4. how do I sum now?
\end{lstlisting}
\item {\footnotesize\color{black!60}2026-11-04}
\begin{lstlisting}[style=prompt]
okay, and how do I square the values of certain axises?
\end{lstlisting}
\item {\footnotesize\color{black!60}2026-11-04}
\begin{lstlisting}[style=prompt]
wait, does sum(axis=(0,2,3)) return 3 numbers?
\end{lstlisting}
\item {\footnotesize\color{black!60}2026-11-04}
\begin{lstlisting}[style=prompt]
how do I elegantly get the values of H and W, multiply them, and keep a running total of num pixels?
\end{lstlisting}
\item {\footnotesize\color{black!60}2026-11-04}
\begin{lstlisting}[style=prompt]
can I unpack a np 1d array of size 3 into 3 different variables using a, b, c = arr?
\end{lstlisting}
\item {\footnotesize\color{black!60}2026-11-04}
\begin{lstlisting}[style=prompt]
review my offline_mean_std, does it look good?
\end{lstlisting}
\item {\footnotesize\color{black!60}2026-11-04}
\begin{lstlisting}[style=prompt]
def offline_mean_std(dataloader: DataLoader):
[... 22 lines omitted ...]
    return [mean_r, mean_g, mean_b], [std_r, std_g, std_b]
\end{lstlisting}
\item {\footnotesize\color{black!60}2026-11-04}
\begin{lstlisting}[style=prompt]
does torch tensor have the same capabilities as numpy?
\end{lstlisting}
\item {\footnotesize\color{black!60}2026-11-04}
\begin{lstlisting}[style=prompt]
does tracking gradients mess up with training at all here if I compute in torch instad of nunmpy?
\end{lstlisting}
\item {\footnotesize\color{black!60}2026-11-04}
\begin{lstlisting}[style=prompt]
how do I make torch.zeros(3) into a float64 value
\end{lstlisting}
\item {\footnotesize\color{black!60}2026-11-04}
\begin{lstlisting}[style=prompt]
def offline_mean_std(dataloader: DataLoader):
[... 22 lines omitted ...]
    return [mean_r, mean_g, mean_b], [std_r, std_g, std_b]
\end{lstlisting}
\item {\footnotesize\color{black!60}2026-11-04}
\begin{lstlisting}[style=prompt]
could I do float64 in the dataset? in __getitem__?
\end{lstlisting}
\item {\footnotesize\color{black!60}2026-11-04}
\begin{lstlisting}[style=prompt]
so how would I do that in offline_mean_std before it gets squared?
\end{lstlisting}
\item {\footnotesize\color{black!60}2026-11-04}
\begin{lstlisting}[style=prompt]
let's start working on the cnn. I could probably just define all the layers in the forward pass right?
\end{lstlisting}
\item {\footnotesize\color{black!60}2026-11-04}
\begin{lstlisting}[style=prompt]
I want to use Conv3d since we are dealing with rgb images. I also don't know the parameters for Conv3d
\end{lstlisting}
\item {\footnotesize\color{black!60}2026-11-04}
\begin{lstlisting}[style=prompt]
out_channels return the number of channels similar to in_channels. kernel size refers to how big the window is, if its 3x3 or something else
\end{lstlisting}
\item {\footnotesize\color{black!60}2026-11-04}
\begin{lstlisting}[style=prompt]
i'm not sure if how out_channels will help me, and if I should make it 16. the second layer in_channel should always match the first layers out_channel. stride does change the output spatial size. i don't think i should add any padding aside from 0 padding when the kernel is out of the frame
\end{lstlisting}
\item {\footnotesize\color{black!60}2026-11-04}
\begin{lstlisting}[style=prompt]
i want to start off with 2 conv layers, a max pooling layer, one fc, and then one softmax. which classes should i use?
\end{lstlisting}
\item {\footnotesize\color{black!60}2026-11-04}
\begin{lstlisting}[style=prompt]
can I reuse the same nn.ReLU object multiple times?
\end{lstlisting}
\item {\footnotesize\color{black!60}2026-11-04}
\begin{lstlisting}[style=prompt]
what should the kernel_size be for 224x224 images
\end{lstlisting}
\item {\footnotesize\color{black!60}2026-11-04}
\begin{lstlisting}[style=prompt]
also should the number of channels reduce per layer?
\end{lstlisting}
\item {\footnotesize\color{black!60}2026-11-04}
\begin{lstlisting}[style=prompt]
what does max pool per layer do?
\end{lstlisting}
\item {\footnotesize\color{black!60}2026-11-04}
\begin{lstlisting}[style=prompt]
I'm pretty sure i do not need to calculate all that. there must be a simpler way to just get the number of inputs for linear without all of that
\end{lstlisting}
\item {\footnotesize\color{black!60}2026-11-04}
\begin{lstlisting}[style=prompt]
is out_features for the fc layer num_classes?
\end{lstlisting}
\item {\footnotesize\color{black!60}2026-11-04}
\begin{lstlisting}[style=prompt]
def __init__(self, num_classes=10):
[... 48 lines omitted ...]
        return self.fc(x)
\end{lstlisting}
\item {\footnotesize\color{black!60}2026-11-04}
\begin{lstlisting}[style=prompt]
what would be the best padding per layer? also how do I use flatten?
\end{lstlisting}
\item {\footnotesize\color{black!60}2026-11-04}
\begin{lstlisting}[style=prompt]
lets use maxpool2d to be the thing that controls downsampling, so that means it'll be after every conv1 + relu layer
\end{lstlisting}
\item {\footnotesize\color{black!60}2026-11-04}
\begin{lstlisting}[style=prompt]
class CNN(nn.Module):
[... 54 lines omitted ...]
 i want to do this
\end{lstlisting}
\item {\footnotesize\color{black!60}2026-11-04}
\begin{lstlisting}[style=prompt]
def forward(self, x):
[... 23 lines omitted ...]
        return self.fc(x)
\end{lstlisting}
\item {\footnotesize\color{black!60}2026-11-04}
\begin{lstlisting}[style=prompt]
def calculate_fc_input():
            dummy = torch.zeros(1, 3, 224, 224)
            out = self.pool(self.conv3(self.pool(self.conv2(self.pool(self.conv1(dummy))))))
            return out.shape
\end{lstlisting}
\item {\footnotesize\color{black!60}2026-11-04}
\begin{lstlisting}[style=prompt]
class CNN(nn.Module):
[... 53 lines omitted ...]
        return self.fc(x)
\end{lstlisting}
\item {\footnotesize\color{black!60}2026-11-04}
\begin{lstlisting}[style=prompt]
let's work on normalization
\end{lstlisting}
\item {\footnotesize\color{black!60}2026-11-04}
\begin{lstlisting}[style=prompt]
I'm going to use this
[... 9 lines omitted ...]
 I also will create a normalize global function
\end{lstlisting}
\item {\footnotesize\color{black!60}2026-11-04}
\begin{lstlisting}[style=prompt]
I will be using only the train mean std. I will hardcode numbers separately. I didn't even think of using combiging ToTensor and normalize in on transform function. Does Normalize(mean, std) take into list/numpy arrays?
\end{lstlisting}
\item {\footnotesize\color{black!60}2026-11-04}
\begin{lstlisting}[style=prompt]
calculate_mean_std contains the mean std for the full dataset, i don't want data leakage
\end{lstlisting}
\item {\footnotesize\color{black!60}2026-11-04}
\begin{lstlisting}[style=prompt]
What I'll do most likely is use the train mean std for local training. How do I build the transform compose function thing
\end{lstlisting}
\item {\footnotesize\color{black!60}2026-11-04}
\begin{lstlisting}[style=prompt]
how do I use a transform?
\end{lstlisting}
\item {\footnotesize\color{black!60}2026-11-04}
\begin{lstlisting}[style=prompt]
it seems like the mean std stats for all 3 aren't that different after all. i'll just use calculate_mean_std
\end{lstlisting}
\item {\footnotesize\color{black!60}2026-11-04}
\begin{lstlisting}[style=prompt]
is there another transform that clips images to 224x224?
\end{lstlisting}
\item {\footnotesize\color{black!60}2026-11-04}
\begin{lstlisting}[style=prompt]
Right after normalize
\end{lstlisting}
\item {\footnotesize\color{black!60}2026-11-04}
\begin{lstlisting}[style=prompt]
can you do resize on a pil even though its from pytorch?
\end{lstlisting}
\item {\footnotesize\color{black!60}2026-11-04}
\begin{lstlisting}[style=prompt]
for training, which kwargs should I be aware of? I know of different optimizers, epochs, etc. I'm going to employ a custom grid search mechanism to find hyperparameters
\end{lstlisting}
\item {\footnotesize\color{black!60}2026-11-04}
\begin{lstlisting}[style=prompt]
i'm using a mac m5. what hsould be the device?
\end{lstlisting}
\item {\footnotesize\color{black!60}2026-11-04}
\begin{lstlisting}[style=prompt]
does using certain devices change the structure of models.pt?
\end{lstlisting}
\item {\footnotesize\color{black!60}2026-11-04}
\begin{lstlisting}[style=prompt]
what is the pseudocode for train and test?
\end{lstlisting}
\item {\footnotesize\color{black!60}2026-11-04}
\begin{lstlisting}[style=prompt]
for the hyperparameters, give me a clean list of which ones I should tune and the range/options of each one
\end{lstlisting}
\item {\footnotesize\color{black!60}2026-11-04}
\begin{lstlisting}[style=prompt]
are the name of the hyperparameters the exact keyword that we pass in? how do I createoptimizer and loss function and scheduler?
\end{lstlisting}
\item {\footnotesize\color{black!60}2026-11-04}
\begin{lstlisting}[style=prompt]
how do I use test function as the val accuracy?
\end{lstlisting}
\item {\footnotesize\color{black!60}2026-11-04}
\begin{lstlisting}[style=prompt]
is model.to(device=device) idempotent?
\end{lstlisting}
\item {\footnotesize\color{black!60}2026-11-04}
\begin{lstlisting}[style=prompt]
how to disable grad tracking, and reneable it
\end{lstlisting}
\item {\footnotesize\color{black!60}2026-11-04}
\begin{lstlisting}[style=prompt]
how do I increment total and correct?
\end{lstlisting}
\item {\footnotesize\color{black!60}2026-11-04}
\begin{lstlisting}[style=prompt]
how should I initialize optim.Adam
\end{lstlisting}
\item {\footnotesize\color{black!60}2026-11-04}
\begin{lstlisting}[style=prompt]
how to zero out optimzer acculated gradients
\end{lstlisting}
\item {\footnotesize\color{black!60}2026-11-04}
\begin{lstlisting}[style=prompt]
based on my grid_search, how do I create the scheduler?
\end{lstlisting}
\item {\footnotesize\color{black!60}2026-11-04}
\begin{lstlisting}[style=prompt]
how do I update running_loss in train?
\end{lstlisting}
\item {\footnotesize\color{black!60}2026-11-04}
\begin{lstlisting}[style=prompt]
this means I should keep a running total of images seen as well so that I can get a good mean
\end{lstlisting}
\item {\footnotesize\color{black!60}2026-11-04}
\begin{lstlisting}[style=prompt]
how do I use rich.progress to log the loops and everything
\end{lstlisting}
\item {\footnotesize\color{black!60}2026-11-04}
\begin{lstlisting}[style=prompt]
who it be batch_task advance=1 or advance=batch_size?
\end{lstlisting}
\item {\footnotesize\color{black!60}2026-11-04}
\begin{lstlisting}[style=prompt]
what's the default for DataLoader classes?
\end{lstlisting}
\item {\footnotesize\color{black!60}2026-11-04}
\begin{lstlisting}[style=prompt]
how do I report running_loss for epoch, curr_val_accuracy, using progress
\end{lstlisting}
\item {\footnotesize\color{black!60}2026-11-04}
\begin{lstlisting}[style=prompt]
how do I log the fact that val_accuracy for a certain epoch is the best and that I want to inform the user that I saved it
\end{lstlisting}
\item {\footnotesize\color{black!60}2026-11-04}
\begin{lstlisting}[style=prompt]
I don't want to save model.pt into a new file. I want to keep it as a variable for now. is that possible?
\end{lstlisting}
\item {\footnotesize\color{black!60}2026-11-04}
\begin{lstlisting}[style=prompt]
is copy part of standard python library?
\end{lstlisting}
\item {\footnotesize\color{black!60}2026-11-04}
\begin{lstlisting}[style=prompt]
def train(model: CNN, train_loader: DataLoader, val_loader: DataLoader, **kwargs):
[... 44 lines omitted ...]
        return best_model_weights, best_val_accuracy
\end{lstlisting}
\item {\footnotesize\color{black!60}2026-11-04}
\begin{lstlisting}[style=prompt]
I want to add a task in progress that says what hyperparamters I'm using, how do I do that?
\end{lstlisting}
\item {\footnotesize\color{black!60}2026-11-04}
\begin{lstlisting}[style=prompt]
read through @hw4_ske.py and review for any errors
\end{lstlisting}
\item {\footnotesize\color{black!60}2026-11-04}
\begin{lstlisting}[style=prompt]
running_loss = 0
            total_images = 0
            for images, labels in train_loader:
                images = images.to(device)
                labels = labels.to(device)
\end{lstlisting}
\item {\footnotesize\color{black!60}2026-11-05}
\begin{lstlisting}[style=prompt]
review @hw4_ske.py
\end{lstlisting}
\item {\footnotesize\color{black!60}2026-11-05}
\begin{lstlisting}[style=prompt]
progress.update(
[... 5 lines omitted ...]
 this line does not change description at all
\end{lstlisting}
\item {\footnotesize\color{black!60}2026-11-05}
\begin{lstlisting}[style=prompt]
with Progress(
[... 43 lines omitted ...]
    return best_model, best_val_accuracy
\end{lstlisting}
\item {\footnotesize\color{black!60}2026-11-05}
\begin{lstlisting}[style=prompt]
for some reason, loss is stuck at 1.3322 and the val_acc is not going up or down
\end{lstlisting}
\item {\footnotesize\color{black!60}2026-11-05}
\begin{lstlisting}[style=prompt]
why don't you check all that for me and tell me what is going on?
\end{lstlisting}
\item {\footnotesize\color{black!60}2026-11-05}
\begin{lstlisting}[style=prompt]
I think I found out the issue. There were too many inputs going into the fully connected layer. I've decided to add new layers into the network, and now the loss converages better. Let's move into writing the report, where I'll put my work in. Follow how the report in @../3/ was generated. It doesn't seem like I need to write up anything for the programming part. Add sections for the Math Questions and the AI usuage part, and let's begin
\end{lstlisting}
\item {\footnotesize\color{black!60}2026-11-05}
\begin{lstlisting}[style=prompt]
change all prompts that have more than a certain number of lines into a summary of what it was so that the page length is shorther
\end{lstlisting}
\end{enumerate}
